\subsection{\texorpdfstring{\(\mathbf{R}^n\) 中的欧氏距离}{R 的 n 次幂中的欧氏距离}}

按照一般定义，两点 \(\boldsymbol{x} = (x_{i})\) 与 \(\boldsymbol{y} = (y_{i})\) 之间的欧氏距离是数

\[
d (\boldsymbol {x}, \boldsymbol {y}) = \sqrt {\sum_ {i = 1} ^ {n} (x _ {i} - y _ {i}) ^ {2}} \geqslant 0.
\]

我们回忆它的主要性质。关系 \(d(x, y) = o\) 等价于 \(x = y\) 。我们有 \(d(x, y) = d(y, x)\) ；对一切标量 \(t \in \mathbf{R}\) ，有 \(d(tx, ty) = |t|d(x, y)\) ；对一切 \(z \in \mathbf{R}^n\) ，有 \(d(x + z, y + z) = d(x, y)\) ；换句话说，两点间的距离在平移下不变。距离 \(d(o, x)\) ——从原点 \(o\) 到一点 \(x\) 的距离——也记作 \(||x||\) ，称为 \(x\) 的欧氏范数（或在不致混淆时简称 \(x\) 的范数；参见第 IX 章，§ 3，第 3 小节）。于是 \(d(x, y) = ||y - x||\) 。

当 \(n = 1\) 时，点 \(x, y\) （\(\mathbf{R}\) 的）之间的欧氏距离归结为长度 \(|y - x|\) （以 \(x\) 与 \(y\) 为端点的区间的长度）。对任意 \(n\) ，我们说 \(d(x, y) = ||y - x||\) 是以 \(x\) 与 \(y\) 为端点的线段的长度。

欧氏距离满足三角不等式

\[
d (\boldsymbol {x}, \boldsymbol {y}) \leqslant d (\boldsymbol {x}, \boldsymbol {z}) + d (\boldsymbol {z}, \boldsymbol {y})\tag{1}
\]

对 \(R^{n}\) 中一切 x, y, z 成立。

我们回忆，(1) 的证明归结为下述不等式的证明：

\[
\left(\sum_ {i = 1} ^ {n} \left(x _ {i} + y _ {i}\right) ^ {2}\right) ^ {\frac {1}{2}} \leqslant \left(\sum_ {i = 1} ^ {n} x _ {i} ^ {2}\right) ^ {\frac {1}{2}} + \left(\sum_ {i = 1} ^ {n} y _ {i} ^ {2}\right) ^ {\frac {1}{2}};
\]

它又等价于柯西–施瓦茨不等式

\[
\left(\sum_ {i = 1} ^ {n} x _ {i} y _ {i}\right) ^ {2} \leqslant \left(\sum_ {i = 1} ^ {n} x _ {i} ^ {2}\right) \left(\sum_ {i = 1} ^ {n} y _ {i} ^ {2}\right),
\]

后者是拉格朗日恒等式的直接推论

\[
\left(\sum_ {i = 1} ^ {n} x _ {i} ^ {2}\right) \left(\sum_ {i = 1} ^ {n} y _ {i} ^ {2}\right) - \left(\sum_ {i = 1} ^ {n} x _ {i} y _ {i}\right) ^ {2} = \frac {1}{2} \sum_ {i, j} (x _ {i} y _ {j} - x _ {j} y _ {i}) ^ {2}.
\]

这个证明同时表明，(1) 的两边仅当 z 是以 x 与 y 为端点的线段的点时才能相等。

由 (1) 我们推出不等式

\[
d (\boldsymbol {x}, \boldsymbol {y}) \geqslant | d (\boldsymbol {x}, \boldsymbol {z}) - d (\boldsymbol {y}, \boldsymbol {z}) |.\tag{2}
\]

最后，若 \(\pmb{x} = (x_{i}),\pmb{y} = (y_{i})\) ，我们有

\[
\sup _ {1 \leqslant i \leqslant n} | x _ {i} - y _ {i} | \leqslant d (x, y) \leqslant \sqrt {n}. \sup _ {1 \leqslant i \leqslant n} | x _ {i} - y _ {i} |.\tag{3}
\]

因此，\(\mathbf{R}^n\) 的子集 A 是有界的（§ 1，第 1 小节）当且仅当

\[
\sup _ {\boldsymbol {x} \in \mathbf {A}} | | \boldsymbol {x} | | <   + \infty .
\]

